% TODO proof-reading
\subsection{beep()の呼び出し}

これは引数なしで2つの関数を呼び出すシンプルな例です：

\begin{lstlisting}[style=customjava]
	public static void main(String[] args)
	{
		java.awt.Toolkit.getDefaultToolkit().beep();
	};
\end{lstlisting}

\begin{lstlisting}
  public static void main(java.lang.String[]);
    flags: ACC_PUBLIC, ACC_STATIC
    Code:
      stack=1, locals=1, args_size=1
         0: invokestatic  #2      // Method java/awt/Toolkit.getDefaultToolkit:()Ljava/awt/Toolkit;
         3: invokevirtual #3      // Method java/awt/Toolkit.beep:()V
         6: return        
\end{lstlisting}

オフセット0の最初の\TT{invokestatic}は\\
\TT{java.awt.Toolkit.getDefaultToolkit()}を呼び出し、
\TT{Toolkit}クラスのオブジェクトへの参照を返します。\\
オフセット3の\TT{invokevirtual}命令は、このクラスの\TT{beep()}メソッドを呼び出します。